O Aprendizado Federado (Federated Learning -- FL) é um paradigma de treinamento distribuído em que diversos clientes colaboram na construção de um modelo global sem compartilhar diretamente seus dados, enviando ao servidor apenas atualizações do modelo. Entretanto, em cenários adversos, clientes maliciosos ou falhos podem manipular essas atualizações (por exemplo, por comportamento bizantino e envenenamento), comprometendo a convergência do treinamento e degradando o desempenho do modelo final. Nesse contexto, a principal dificuldade consiste em definir mecanismos de agregação que reduzam a influência de atualizações anômalas, mantendo desempenho competitivo em condições realistas. Este trabalho tem como objetivo avaliar algoritmos de agregação robusta aplicados ao FL sob diferentes níveis de adversarialidade. Para isso, serão comparados métodos robustos de agregação, como Krum e variantes baseadas em estatísticas robustas (por exemplo, mediana coordenada-a-coordenada e média aparada), frente a um agregador de referência baseado em média. O desempenho será analisado por métricas de aprendizado (como acurácia e F1-score), estabilidade do treinamento e sensibilidade à proporção de participantes maliciosos, buscando evidenciar vantagens, limitações e condições de aplicabilidade de cada agregador. Espera-se que os resultados auxiliem na escolha de estratégias de agregação robusta para sistemas de FL expostos a participantes não confiáveis.

\palavraschave{aprendizado federado; agregação robusta; falhas bizantinas; envenenamento de modelo; segurança em aprendizado de máquina}



%A Doença de Alzheimer é uma condição neurodegenerativa sem cura que se apresenta como um dos principais desafios em saúde pública mundial, afetando não apenas o paciente, mas também as pessoas ao seu redor. Nesse contexto, a principal dificuldade em relação à doença é a sua detecção precoce, visto que vários de seus sintomas podem ser confundidos com o envelhecimento normal. Este trabalho tem como objetivo validar sinais de áudio baseados na fala como um biomarcador promissor para realizar essa detecção. Para tal, serão treinados e validados dois modelos classificadores — Árvore de Decisão e XGBoost — com foco na utilização da técnica SHAP para promover não apenas a acurácia, mas também a explicabilidade dos modelos. O desempenho será comparado com base em métricas como acurácia, sensibilidade, especificidade, F1-Score e matriz de confusão, bem como na capacidade de explicar quais atributos colaboraram para a decisão final. Espera-se que os resultados obtidos destaquem as vantagens e limitações de cada abordagem, oferecendo \textit{insights} sobre a aplicabilidade de modelos explicáveis no contexto clínico e fornecendo, assim, um melhor suporte ao diagnóstico da Doença de Alzheimer. 

% Separe as palavras-chave por ponto
%\palavraschave{Doença de Alzheimer; Análise de Fala; Aprendizado de Máquina; IA Explicável; Árvore de Decisão; XGBoost.}